% Аннотация (подключается в annotation.tex и main.tex)
Первые слои больших языковых моделей (LLM) на основе архитектуры
Transformer~\cite{vaswani2017attention} выполняют сравнительно простые
операции: восстанавливают статистику биграмм и коротких
$n$-грамм~\cite{elhage2021framework}, реализуют фиксированные позиционные
паттерны внимания и хранят поверхностные шаблоны в слоях
прямого распространения~\cite{geva2021ffn}.
Мы исследуем гипотезу о том, что вычисления первых $k$ слоёв можно
заменить \emph{детерминированными} алгоритмами, то есть процедурами без
случайной выборки и без дообучения градиентным спуском, которые
строятся по весам модели и статистике калибровочного корпуса в
замкнутой форме. Мы адаптируем к каузальному декодеру Qwen3 методы
эффективного внимания и сжатия слоёв MLP из литературы и предлагаем
собственные: факторизацию внимания <<позиционное ядро $\times$
значимость ключевого токена>>, типизированную замену голов,
линеаризацию MLP с компенсацией дрейфа остаточного потока и табличные
замены блоков. Эксперименты проводятся на Qwen3-0.6B и 1.7B
(Base) с оценкой по KL-дивергенции к исходной модели, perplexity и
стандартным бенчмаркам.
Профилирование показывает, что первые слои образуют отдельный режим
<<до стока внимания>>: массивная активация первого токена и сток
внимания возникают лишь после нескольких первых слоёв, что задаёт
естественную глубину замены $k_{\mathrm{sink}}$. Мы также показываем,
что рандомизированные признаки FAVOR+ (Performer~\cite{choromanski2021performer})
неприменимы к первым слоям без переобучения: нормы ключей после
QK-нормализации лежат далеко за пределами области, где их оценка
ядра устойчива.
Основные результаты. (1) Внимание первых трёх слоёв заменяется окном
из 64 позиций со стоком с ростом perplexity 1.5--1.8\,\% и без потерь на
семи бенчмарках; MLP первых слоёв заменяется хуже: лучший метод "---
предложенный детерминированный предиктор активных нейронов "--- стоит
4.7\,\% (0.6B) и 1.7\,\% (1.7B) при четверти вычислений MLP, а по карте
заменяемости MLP ранних слоёв оказываются самыми незаменимыми в модели.
(2) Мы обнаруживаем ловушку калибровки в замкнутой форме: единичные
токены с массивной активацией доминируют в задаче наименьших квадратов,
ошибка ложится в каналы стока внимания и обрушивает модель; мы даём
формулу влияния строки, воспроизводим эффект дозой отравленных окон и
устраняем его исключением строк-выбросов, а также показываем, что
массивную активацию создают один--пять нейронов, удаление которых
разрушает модель. (3) Априорные статистики слоя предсказывают KL его
замены с ранговой корреляцией 0.9--0.97 и позволяют планировать замены
на непросмотренной модели в 10--15 раз точнее замены префикса.
(4) На втором семействе (SmolLM2-1.7B, Llama-3.2-1B) сток оказывается
привязан не к позиции, а к токену: один нейрон создаёт массивную
активацию на <<the>> в любой позиции, и замены переносятся лишь с
защитой таких нейронов и хранением таких ключей (окно: $64\to4\,\%$ PPL
при $k=8$); MLP первых слоёв у этого семейства заменяется хуже, а
планировщик переносится.
Результаты дополняют дискуссию о том, насколько внутренние вычисления
обученных сетей <<уже чисты>> и поддаются явному
описанию~\cite{cammarata_clean,gao_notclean}.

\medskip
\noindent\textbf{Ключевые слова:} трансформеры, большие языковые модели,
аппроксимация внимания, механистическая интерпретируемость,
детерминированные алгоритмы, низкоранговые приближения, сжатие MLP.
